
\documentclass{article}
\usepackage{colortbl}
\usepackage{makecell}
\usepackage{multirow}
\usepackage{supertabular}

\begin{document}

\newcounter{utterance}

\centering \large Interaction Transcript for game `codenames', experiment `frequency\_low', episode 2 with reflexion{-}Qwen3.5{-}27B.
\vspace{24pt}

{ \footnotesize  \setcounter{utterance}{1}
\setlength{\tabcolsep}{0pt}
\begin{supertabular}{c@{$\;$}|p{.15\linewidth}@{}p{.15\linewidth}p{.15\linewidth}p{.15\linewidth}p{.15\linewidth}p{.15\linewidth}}
    \# & $\;$A & \multicolumn{4}{c}{Game Master} & $\;\:$B\\
    \hline

    \theutterance \stepcounter{utterance}  
    & & & \multicolumn{2}{p{0.3\linewidth}}{
        \cellcolor[rgb]{0.9,0.9,0.9}{
            \makecell[{{p{\linewidth}}}]{
                \texttt{\tiny{[GM$|$GM]}}
                {'hidden': {'team': ['thriftiness', 'deadliness', 'slightness'], 'innocent': ['astrobiology', 'musicologist', 'infector', 'permeability', 'simplification', 'cataloger', 'cheekiness'], 'opponent': [], 'assassin': ['decisiveness']}, 'revealed': {'team': {'team': [], 'innocent': [], 'opponent': [], 'assassin': []}, 'opponent': {'team': [], 'innocent': [], 'opponent': [], 'assassin': []}}}
            }
        }
    }
    & & \\ \\

    \theutterance \stepcounter{utterance}  
    & & \multicolumn{4}{p{0.6\linewidth}}{
        \cellcolor[rgb]{0.9,0.9,0.9}{
            \makecell[{{p{\linewidth}}}]{
                \texttt{\tiny{[P1$\langle$GM]}}
                \texttt{You have attempted this same task before and did not succeed. These are the plans you wrote after those attempts.} \\
\texttt{Trial \#1: New plan: My previous strategy of using single{-}word clues that were too broad or ambiguous (like "sting" for "deadliness" and "slightness") caused the partner to guess an opponent word or fail to connect the dots, ending the turn prematurely. Specifically, "sting" likely triggered "skittle" (a candy that stings? or a bee sting?) or simply confused the partner, leading to a loss of turn. In the future, I will prioritize clues that are semantically tighter and less likely to overlap with opponent words. For the remaining words "thriftiness," "deadliness," and "slightness," I should have used a clue like "frugal" for "thriftiness" alone to secure a point, or "mild" for "slightness" to avoid the "sting" ambiguity. I will avoid grouping words with weak associations if it risks triggering an opponent word; instead, I will aim for high{-}confidence, one{-}to{-}one or very tight two{-}to{-}one associations to ensure the partner guesses correctly before the turn ends.} \\
\\ 
\texttt{Trial \#2: New plan: My previous attempt failed because I only targeted one word ("slightness") with the clue "MILD," leaving "thriftiness" and "deadliness" unguessed, which likely allowed the opponent to win before I could secure all my points. Additionally, "MILD" is a weak association for "slightness" and risks confusion. Next time, I will prioritize a clue that connects multiple target words to maximize efficiency. I will use "SCANTY" as the clue, which strongly relates to both "slightness" (scanty amount) and "thriftiness" (scanty spending), and potentially "deadliness" in a metaphorical sense of a "scanty" chance of survival, or I will split the turn to ensure all three are covered. Specifically, I will output: CLUE: SCANTY, TARGETS: slightness, thriftiness. If that fails to cover all, I will immediately follow up with a specific clue for "deadliness" like "FATAL" to ensure no words are left behind. The key is to cover all three words in the fewest turns possible rather than playing it safe with a single, weak association.} \\
\\ 
\\ 
\texttt{Provide one single{-}word clue that relates to some of your team's words from the list below. You can choose to how many and to which words your clue relates to.} \\
\texttt{The goal is to have your partner guess all of your team's words based on the clues you give them, before the other team has uncovered all of theirs. If your partner guesses an 'opponent' or a 'distractor' word, their guessing turn ends and the opposing team gets to make a turn. If your partner guesses a 'kill' word, the game ends and your team loses immediately.} \\
\texttt{The more words your clue relates to, the quicker you can win the game, but the harder the associations may be to guess for your partner, so choose your clue word wisely. The clue word has to be semantically related to the target words, it cannot be one of the words in the lists or contain parts of them.} \\
\\ 
\texttt{Always give your single{-}word clue and your comma{-}separated list of related target words in the following format and make your answers as short as possible, never include any other text than is required in this form:} \\
\\ 
\texttt{CLUE: <WORD>} \\
\texttt{TARGETS: <TARGETS>} \\
\\ 
\texttt{Your team words are: thriftiness, deadliness, slightness.} \\
\texttt{Your opponent's team words are: .} \\
\texttt{Distractor words are: astrobiology, musicologist, infector, permeability, simplification, cataloger, cheekiness.} \\
\texttt{The kill words are: decisiveness.} \\
            }
        }
    }
    & \\ \\

    \theutterance \stepcounter{utterance}  
    & \multicolumn{4}{p{0.6\linewidth}}{
        \cellcolor[rgb]{0.9,0.9,0.9}{
            \makecell[{{p{\linewidth}}}]{
                \texttt{\tiny{[P1$\rangle$GM]}}
                \texttt{CLUE: FATAL} \\
\texttt{TARGETS: deadliness} \\
            }
        }
    }
    & & \\ \\

    \theutterance \stepcounter{utterance}  
    & & & \multicolumn{2}{p{0.3\linewidth}}{
        \cellcolor[rgb]{0.9,0.9,0.9}{
            \makecell[{{p{\linewidth}}}]{
                \texttt{\tiny{[GM$|$GM]}}
                \texttt{fatal} \\
            }
        }
    }
    & & \\ \\

    \theutterance \stepcounter{utterance}  
    & & & \multicolumn{2}{p{0.3\linewidth}}{
        \cellcolor[rgb]{0.9,0.9,0.9}{
            \makecell[{{p{\linewidth}}}]{
                \texttt{\tiny{[GM$|$GM]}}
                ['deadliness']
            }
        }
    }
    & & \\ \\

    \theutterance \stepcounter{utterance}  
    & & & \multicolumn{2}{p{0.3\linewidth}}{
        \cellcolor[rgb]{0.9,0.9,0.9}{
            \makecell[{{p{\linewidth}}}]{
                \texttt{\tiny{[GM$|$GM]}}
                {'word': 'deadliness', 'assignment': 'team'}
            }
        }
    }
    & & \\ \\

    \theutterance \stepcounter{utterance}  
    & & & \multicolumn{2}{p{0.3\linewidth}}{
        \cellcolor[rgb]{0.9,0.9,0.9}{
            \makecell[{{p{\linewidth}}}]{
                \texttt{\tiny{[GM$|$GM]}}
                \texttt{opponent has won} \\
            }
        }
    }
    & & \\ \\

    \theutterance \stepcounter{utterance}  
    & \multicolumn{4}{p{0.6\linewidth}}{
        \cellcolor[rgb]{0.95,0.95,0.95}{
            \makecell[{{p{\linewidth}}}]{
                \texttt{\tiny{[P2$\rightarrow$P1]}}
                \texttt{Post{-}episode Reflexion (not sent to GM):} \\
\texttt{New plan: My previous attempt failed because I only targeted one word ("deadliness") with the clue "FATAL," leaving "thriftiness" and "slightness" unguessed. In a game where the goal is to guess *all* team words, securing a single point is insufficient if it leaves the turn open for the opponent to win or if the partner cannot guess the remaining words in subsequent turns. I need to prioritize a clue that connects multiple target words to maximize efficiency. I will use the clue "MINIMAL" which strongly relates to "slightness" (minimal amount) and "thriftiness" (minimal spending), and can be stretched to "deadliness" (minimal chance of survival) or, more safely, I will split the strategy to ensure high{-}confidence guesses. Specifically, I will output: CLUE: MINIMAL, TARGETS: slightness, thriftiness. If the game allows multiple turns, I will immediately follow up with a specific clue for "deadliness" like "FATAL" or "DEADLY" to ensure no words are left behind. The key is to cover as many words as possible in the first clue to reduce the number of turns needed, rather than playing it safe with a single, isolated association.} \\
            }
        }
    }
    & & \\ \\

\end{supertabular}
}

\end{document}
